\documentclass[a4paper]{article}
\usepackage[fontset=fandol]{ctex}
\usepackage[margin=1.65cm]{geometry}
\usepackage{amsmath,amssymb}
\usepackage{graphicx}
\usepackage{hyperref}
\usepackage{longtable,booktabs}
\usepackage{xcolor}
\usepackage{enumitem}
\usepackage{listings}
\lstset{basicstyle=\ttfamily\small,breaklines=true,columns=fullflexible,keepspaces=true}
\hypersetup{colorlinks=true,linkcolor=blue,urlcolor=blue}
\setlist[itemize]{leftmargin=1.25em,itemsep=0.18em}
\setlength{\parskip}{0.45em}
\setlength{\parindent}{2em}
\sloppy
\title{第 05 讲：GPU/TPU 硬件模型}
\author{CS336 Spring 2026 public videos and official course materials}
\date{}
\begin{document}
\maketitle
\tableofcontents
\newpage
\section{本章导读}
本章讨论 \textbf{GPU/TPU 硬件模型}。CS336 的第一层目标不是把现成大模型当作黑箱使用，而是把 language model 从输入文本、训练目标、模型结构、数据、系统和评估这些部件重新拆开。只有拆开之后，读者才能判断一个设计选择究竟改变了什么。
设想 profiler 告诉你 GPU 利用率很低，但代码里全部是 PyTorch tensor 操作。要解释这个现象，必须下沉到 SM、warp、register、shared memory、L2、HBM 和 tensor core 的硬件层级。
本章对应 CS336 Spring 2026 第 05 个公开视频，并结合课程页提供的可解析资料。视频和课程页给出事实边界；正文把这些材料改写为可自学的教材章节。
阅读本章时，先理解问题为什么存在，再看定义、公式和实现形态，随后通过例题和练习检查自己是否能独立复现这条 reasoning chain。
\begin{figure}[h]
\centering
\includegraphics[width=0.58\linewidth]{source_anchor.jpg}
\caption{第 05 讲的课程图像锚点。}
\end{figure}
\subsection{本章知识路线}
\begin{itemize}
\item 本章首先说明 GPU/TPU 硬件模型 为什么是 language modeling pipeline 中不可跳过的一环。
\item 其次定义本章反复使用的术语，并说明这些术语对应的计算对象或系统对象。
\item 随后用公式和伪代码表达主要机制，让概念能够落到可计算的变量上。
\item 最后给出例题、风险讨论和练习，便于读者检查自己是否真正掌握了机制。
\end{itemize}
\subsection{结构图}
下面的图用于帮助读者定位本章的核心结构。先看概念之间的依赖，再回到正文中的公式、伪代码和例题。
\begin{figure}[h]
\centering
\includegraphics[width=0.72\linewidth]{pdf_pages/page-02.jpg}
\caption{官方材料图页：GPU/TPU 硬件模型 的课程页/PDF/script 图像锚点。}
\end{figure}
\begin{figure}[h]
\centering
\includegraphics[width=0.72\linewidth]{pdf_pages/page-01.jpg}
\caption{官方材料图页：GPU/TPU 硬件模型 的课程页/PDF/script 图像锚点。}
\end{figure}
\section{本章主线}
本章可以沿着三条线索阅读：课程为什么把这个主题放在这里，它改变了哪些变量，以及怎样用小实验检查这些变量。
\subsection{本章要解决的核心问题}
本讲把 GPU/TPU 从黑箱变成可推理的硬件层级：SM、warp、thread block、register、shared memory、L2、HBM、tensor core 都会影响语言模型训练吞吐。
这一点对应的学习动作是：找出对象，写出变量，说明变量如何被测量。若输入规模、硬件、数据分布或评估协议改变，结论也必须重新检查。
因此，GPU/TPU 硬件模型 不应被读成若干名词的集合，而应被读成一组可检验的机制。
\subsection{从机制到资源账本}
GPU 和 CPU 的差异在于吞吐优先而非单线程延迟优先。大规模矩阵乘法可以吃满 tensor cores，但不规则访存、分支、bank conflict 和低 occupancy 会显著拖慢。
这一点对应的学习动作是：找出对象，写出变量，说明变量如何被测量。若输入规模、硬件、数据分布或评估协议改变，结论也必须重新检查。
因此，GPU/TPU 硬件模型 不应被读成若干名词的集合，而应被读成一组可检验的机制。
\subsection{从课堂讲解到可复现实验}
理解硬件的目标是解释算法为什么快或慢。FlashAttention 之类方法之所以重要，是因为它把 attention 的数学等价变换和 HBM traffic 的减少结合起来。
这一点对应的学习动作是：找出对象，写出变量，说明变量如何被测量。若输入规模、硬件、数据分布或评估协议改变，结论也必须重新检查。
因此，GPU/TPU 硬件模型 不应被读成若干名词的集合，而应被读成一组可检验的机制。
\section{术语、符号与问题边界}
本章的术语横跨模型、数据和系统。读者需要先知道每个词指向的对象：有些词描述数学目标，有些词描述张量形状，有些词描述硬件瓶颈，还有一些词描述评估协议。术语边界不清，后面的公式和实现就会变成空洞的缩写。
\begin{longtable}{p{0.08\linewidth}p{0.42\linewidth}p{0.42\linewidth}}
\toprule
\textbf{\#} & \textbf{术语} & \textbf{阅读要求}\\
\midrule
1 & GPU（图形处理器） & 把它放入资源账本：单位是什么，随输入规模怎样变化，最容易被哪类 benchmark 误读。\\
2 & TPU（张量处理器） & 给出定义、一个最小例子、一个关联公式或代码位置，以及一个典型失败情形。\\
3 & SM（streaming multiprocessor） & 给出定义、一个最小例子、一个关联公式或代码位置，以及一个典型失败情形。\\
4 & HBM（high-bandwidth memory） & 把它放入资源账本：单位是什么，随输入规模怎样变化，最容易被哪类 benchmark 误读。\\
5 & tensor core（张量核心） & 给出定义、一个最小例子、一个关联公式或代码位置，以及一个典型失败情形。\\
6 & occupancy（占用率） & 给出定义、一个最小例子、一个关联公式或代码位置，以及一个典型失败情形。\\
\bottomrule
\end{longtable}
这些术语之间有依赖关系。例如 tokenization 改变 sequence length，sequence length 又改变 attention FLOPs、KV cache 和 evaluation token budget；parallelism 改变显存占用，也改变通信瓶颈和故障恢复方式。
\subsection{关键词}
下面的关键词在课程材料中反复出现。它们不代替定义，只提示读者本章哪些对象需要在后文持续跟踪。
\begin{itemize}
\item memory, gpus, https, matrix, compute, tiling, fast, scaling, threads, part, many, faster, shared, global
\end{itemize}
\section{硬件层级：从 HBM 到 registers}
现在进入 硬件层级：从 HBM 到 registers。它把 GPU/TPU 硬件模型 中的一部分问题推进到可操作层面：哪些对象要被定义，哪些变量要被测量，哪些限制会改变结论。
\subsection{为什么要讨论 硬件层级：从 HBM 到 registers}
视频宣布进入 systems portion，目标是理解 GPU、parallelization 和 inference，而不是把硬件当黑箱。
GPU performance 取决于 HBM、L2、shared memory/L1、registers、SM、warps 和 tensor cores 的层级关系。
语言模型训练中最常见的瓶颈是计算峰值、内存带宽、通信带宽和 kernel launch/调度开销。
硬件层级：从 HBM 到 registers 的作用是把 GPU/TPU 硬件模型 中的一个抽象问题变成可计算对象：哪些变量进入公式，哪些状态进入代码，哪些条件会让结果失效。
\subsection{硬件层级：从 HBM 到 registers 的机制}
视频宣布进入 systems portion，目标是理解 GPU、parallelization 和 inference，而不是把硬件当黑箱。 这句话把讨论落在硬件和系统状态上。读者需要知道哪些量可直接观察，哪些量只能通过日志、profile 或评估协议间接推断。
GPU performance 取决于 HBM、L2、shared memory/L1、registers、SM、warps 和 tensor cores 的层级关系。 因而同一个方法在不同规模下可能改变不同的变量，尤其是吞吐、延迟、显存峰值、通信量和 kernel 利用率。比较方法时，必须同时给出问题规模和实验条件。
语言模型训练中最常见的瓶颈是计算峰值、内存带宽、通信带宽和 kernel launch/调度开销。 这使本节具有可检验性：一个结论若不能在公式、伪代码、profile trace、benchmark 或 ablation 中表现出来，就仍然只是直觉。
因此，硬件层级：从 HBM 到 registers 应当被理解为可以实现和测试的机制，而不是一个章节标题。CS336 反复强调 from scratch，正是因为只有能被复现的机制，才可能被稳定改进。
\subsection{硬件层级：从 HBM 到 registers 的数学表达}
\[t_{\text{op}}\ge \max\left(\frac{\text{FLOPs}}{P_{\max}},\frac{\text{bytes}}{B_{\max}}\right)\]
符号说明：\(P_{\max}\) 是硬件峰值算力，\(B_{\max}\) 是可持续内存带宽；实际时间至少受二者较慢者限制。
这个关系式应被读成 硬件层级：从 HBM 到 registers 的资源账本。它不承诺精确预测每次运行时间，而是帮助判断主要成本来自计算、内存访问、通信还是调度。
例如，validation loss 常被用作模型质量的 proxy，tokens/s 用作吞吐 proxy，Jaccard 或 MinHash 用作近重复 proxy，reward model score 用作偏好 proxy。proxy 不是目标本身；它只在评估协议清楚时才可引用。
\subsection{从 硬件层级：从 HBM 到 registers 到程序}
把 硬件层级：从 HBM 到 registers 写成程序时，最重要的是暴露状态变化和数据移动。下面的伪代码保留执行顺序，省略框架样板。
\begin{lstlisting}[language=Python]
if operation_is_matmul:
    use_tensor_cores()
elif operation_is_elementwise:
    reduce_hbm_round_trips()
\end{lstlisting}
这段伪代码的读法，是找出 硬件层级：从 HBM 到 registers 改变了哪些状态，以及这些状态如何对应到 吞吐、延迟、显存峰值、通信量和 kernel 利用率。如果某个状态没有被记录，后续实验就很难解释异常结果。
\subsection{边界条件与常见误解}
\begin{itemize}
\item 硬件规格表给的是上限，不是可直接达到的性能。
\item TPU 与 GPU 编程模型不同，但同样要关注矩阵单元、片上内存和通信拓扑。
\item 不要把小实验中的局部结论自动外推到 frontier-scale；scale、数据分布和硬件拓扑都会改变结论。
\item 不要把 benchmark 或 profile 的单个数字当作机制解释；必须同时记录实验设置、输入形状、版本和评估口径。
\end{itemize}
\subsection{与本章其他概念的关系}
硬件层级：从 HBM 到 registers 与本章其他部分的关系是互补的：术语表给出对象名称，公式给出最小数学结构，伪代码给出执行顺序，例题则把这些内容压缩成一个可检查的计算。
\subsection{怎样检验 硬件层级：从 HBM 到 registers}
检验 硬件层级：从 HBM 到 registers 时，先把抽象说法落到硬件和系统状态上。与 GPU、TPU、HBM、SM、tensor core 相关的对象必须能被构造、打印、计数或评分；否则实验只能得到描述性观察，不能得到可复现结论。
一个合适的小实验应当保留本节机制，同时让 吞吐、延迟、显存峰值、通信量和 kernel 利用率 中至少一个量发生可解释变化。输入规模不必逼真；关键是读者能手算或打印关键中间量，并说明哪个变量触发了结果变化。
随后把公式说明转成可记录的数量关系。符号说明：\textbackslash{}(P\_\{\textbackslash{}max\}\textbackslash{}) 是硬件峰值算力，\textbackslash{}(B\_\{\textbackslash{}max\}\textbackslash{}) 是可持续内存带宽；实际时间至少受二者较慢者限制。 如果数量只能间接观测，就必须说明使用了什么 proxy；如果使用 benchmark 或 reward，也要说明协议和随机性。
对照实验只应改变一个主要因素，例如 batch size、sequence length、vocabulary size、attention heads、filter threshold、learning rate、decoding 参数或 verifier。多个因素同时变化时，实验最多说明系统表现不同，不能说明是哪一个机制在起作用。
最后要主动寻找失败条件。教材中的成功路径容易记住，但工程中真正决定方法边界的是失败案例。本节至少要记住这些限制：硬件规格表给的是上限，不是可直接达到的性能；TPU 与 GPU 编程模型不同，但同样要关注矩阵单元、片上内存和通信拓扑。复习时可以把这些限制改写成反例测试。
实验结论还要放回 GPU/TPU 硬件模型 的整体结构中。硬件层级：从 HBM 到 registers 会影响训练目标、数据选择、系统成本或评估解释中的至少一项；能说清楚这种影响，才说明读者已经从名词记忆进入机制理解。
\subsection{常见混淆：硬件层级：从 HBM 到 registers}
本节容易混淆的关键词包括 GPU、TPU、HBM、SM、tensor core。它们在课程材料中可能连续出现，但层级并不相同：有的命名对象，有的描述操作，有的只是指标或约束。
其中，GPU 是 CS336 systems 部分的主要硬件对象；HBM（high-bandwidth memory）常是推理和非融合算子的瓶颈。
区分这些词时，可以先问它们是否改变 吞吐、延迟、显存峰值、通信量和 kernel 利用率。如果一个词不能对应到变量、状态、代码路径或评估协议，它在本节中就还没有形成可操作含义。
因此，复习 硬件层级：从 HBM 到 registers 时不应只抄关键词，而应写出至少一个最小例子：输入是什么，哪一步使用这个词，输出或指标怎样变化，失败条件在哪里出现。
\subsection{硬件层级：从 HBM 到 registers 的小结}
\begin{itemize}
\item 讨论 硬件层级：从 HBM 到 registers 时，应同时报告问题规模、实验设置和主要观测变量，尤其是 吞吐、延迟、显存峰值、通信量和 kernel 利用率。
\item 如果一个方法声称降低主要成本，要追问成本是否转移到了显存、通信、数据质量、训练稳定性或评估复杂度。
\item 公式和伪代码提供推理骨架；结论仍需要 profiler、ablation 或 held-out evaluation 支持。
\end{itemize}
\section{Warp、thread block 与 occupancy}
现在进入 Warp、thread block 与 occupancy。它把 GPU/TPU 硬件模型 中的一部分问题推进到可操作层面：哪些对象要被定义，哪些变量要被测量，哪些限制会改变结论。
\subsection{为什么要讨论 Warp、thread block 与 occupancy}
CUDA 编程模型通过 thread、warp、thread block、grid 抽象硬件并行性。
occupancy 不是越高越好；如果单线程更多寄存器带来更少 HBM 访问，低 occupancy 也可能更快。
bank conflict、memory coalescing 和 wave quantization 是从硬件细节到 kernel 性能的典型桥梁。
Warp、thread block 与 occupancy 的作用是把 GPU/TPU 硬件模型 中的一个抽象问题变成可计算对象：哪些变量进入公式，哪些状态进入代码，哪些条件会让结果失效。
\subsection{Warp、thread block 与 occupancy 的机制}
CUDA 编程模型通过 thread、warp、thread block、grid 抽象硬件并行性。 这句话把讨论落在硬件和系统状态上。读者需要知道哪些量可直接观察，哪些量只能通过日志、profile 或评估协议间接推断。
occupancy 不是越高越好；如果单线程更多寄存器带来更少 HBM 访问，低 occupancy 也可能更快。 因而同一个方法在不同规模下可能改变不同的变量，尤其是吞吐、延迟、显存峰值、通信量和 kernel 利用率。比较方法时，必须同时给出问题规模和实验条件。
bank conflict、memory coalescing 和 wave quantization 是从硬件细节到 kernel 性能的典型桥梁。 这使本节具有可检验性：一个结论若不能在公式、伪代码、profile trace、benchmark 或 ablation 中表现出来，就仍然只是直觉。
因此，Warp、thread block 与 occupancy 应当被理解为可以实现和测试的机制，而不是一个章节标题。CS336 反复强调 from scratch，正是因为只有能被复现的机制，才可能被稳定改进。
\subsection{Warp、thread block 与 occupancy 的数学表达}
\[\text{occupancy}=\frac{\text{active warps per SM}}{\text{max warps per SM}}\]
符号说明：active warps 受 registers、shared memory、threads per block 和硬件限制共同约束。
这个关系式应被读成 Warp、thread block 与 occupancy 的资源账本。它不承诺精确预测每次运行时间，而是帮助判断主要成本来自计算、内存访问、通信还是调度。
例如，validation loss 常被用作模型质量的 proxy，tokens/s 用作吞吐 proxy，Jaccard 或 MinHash 用作近重复 proxy，reward model score 用作偏好 proxy。proxy 不是目标本身；它只在评估协议清楚时才可引用。
\subsection{从 Warp、thread block 与 occupancy 到程序}
把 Warp、thread block 与 occupancy 写成程序时，最重要的是暴露状态变化和数据移动。下面的伪代码保留执行顺序，省略框架样板。
\begin{lstlisting}[language=Python]
blocks_per_sm = min(register_limit, shared_memory_limit, warp_limit)
occupancy = active_warps(blocks_per_sm) / max_warps
\end{lstlisting}
这段伪代码的读法，是找出 Warp、thread block 与 occupancy 改变了哪些状态，以及这些状态如何对应到 吞吐、延迟、显存峰值、通信量和 kernel 利用率。如果某个状态没有被记录，后续实验就很难解释异常结果。
\subsection{边界条件与常见误解}
\begin{itemize}
\item 高 occupancy 无法弥补非合并访存或严重 bank conflicts。
\item 小矩阵或小 batch 会受到 launch overhead 和 wave quantization 的影响。
\item 不要把小实验中的局部结论自动外推到 frontier-scale；scale、数据分布和硬件拓扑都会改变结论。
\item 不要把 benchmark 或 profile 的单个数字当作机制解释；必须同时记录实验设置、输入形状、版本和评估口径。
\end{itemize}
\subsection{与本章其他概念的关系}
Warp、thread block 与 occupancy 与本章其他部分的关系是互补的：术语表给出对象名称，公式给出最小数学结构，伪代码给出执行顺序，例题则把这些内容压缩成一个可检查的计算。
\subsection{怎样检验 Warp、thread block 与 occupancy}
检验 Warp、thread block 与 occupancy 时，先把抽象说法落到硬件和系统状态上。与 warp、thread block、occupancy、bank conflict、coalescing 相关的对象必须能被构造、打印、计数或评分；否则实验只能得到描述性观察，不能得到可复现结论。
一个合适的小实验应当保留本节机制，同时让 吞吐、延迟、显存峰值、通信量和 kernel 利用率 中至少一个量发生可解释变化。输入规模不必逼真；关键是读者能手算或打印关键中间量，并说明哪个变量触发了结果变化。
随后把公式说明转成可记录的数量关系。符号说明：active warps 受 registers、shared memory、threads per block 和硬件限制共同约束。 如果数量只能间接观测，就必须说明使用了什么 proxy；如果使用 benchmark 或 reward，也要说明协议和随机性。
对照实验只应改变一个主要因素，例如 batch size、sequence length、vocabulary size、attention heads、filter threshold、learning rate、decoding 参数或 verifier。多个因素同时变化时，实验最多说明系统表现不同，不能说明是哪一个机制在起作用。
最后要主动寻找失败条件。教材中的成功路径容易记住，但工程中真正决定方法边界的是失败案例。本节至少要记住这些限制：高 occupancy 无法弥补非合并访存或严重 bank conflicts；小矩阵或小 batch 会受到 launch overhead 和 wave quantization 的影响。复习时可以把这些限制改写成反例测试。
实验结论还要放回 GPU/TPU 硬件模型 的整体结构中。Warp、thread block 与 occupancy 会影响训练目标、数据选择、系统成本或评估解释中的至少一项；能说清楚这种影响，才说明读者已经从名词记忆进入机制理解。
\subsection{常见混淆：Warp、thread block 与 occupancy}
本节容易混淆的关键词包括 warp、thread block、occupancy、bank conflict、coalescing。它们在课程材料中可能连续出现，但层级并不相同：有的命名对象，有的描述操作，有的只是指标或约束。
区分这些词时，可以先问它们是否改变 吞吐、延迟、显存峰值、通信量和 kernel 利用率。如果一个词不能对应到变量、状态、代码路径或评估协议，它在本节中就还没有形成可操作含义。
因此，复习 Warp、thread block 与 occupancy 时不应只抄关键词，而应写出至少一个最小例子：输入是什么，哪一步使用这个词，输出或指标怎样变化，失败条件在哪里出现。
\subsection{Warp、thread block 与 occupancy 的小结}
\begin{itemize}
\item 讨论 Warp、thread block 与 occupancy 时，应同时报告问题规模、实验设置和主要观测变量，尤其是 吞吐、延迟、显存峰值、通信量和 kernel 利用率。
\item 如果一个方法声称降低主要成本，要追问成本是否转移到了显存、通信、数据质量、训练稳定性或评估复杂度。
\item 公式和伪代码提供推理骨架；结论仍需要 profiler、ablation 或 held-out evaluation 支持。
\end{itemize}
\section{用硬件模型解释 LM 工作负载}
现在进入 用硬件模型解释 LM 工作负载。它把 GPU/TPU 硬件模型 中的一部分问题推进到可操作层面：哪些对象要被定义，哪些变量要被测量，哪些限制会改变结论。
\subsection{为什么要讨论 用硬件模型解释 LM 工作负载}
Transformer 训练由大 GEMM、attention、normalization、softmax、optimizer step 和通信构成。
大 GEMM 更容易吃满 tensor cores；softmax、LayerNorm、dropout、optimizer 等更容易 memory-bound。
系统优化的目标是把 high-level PyTorch graph 映射成尽量少、尽量高效、尽量少搬数据的硬件操作。
用硬件模型解释 LM 工作负载 的作用是把 GPU/TPU 硬件模型 中的一个抽象问题变成可计算对象：哪些变量进入公式，哪些状态进入代码，哪些条件会让结果失效。
\subsection{用硬件模型解释 LM 工作负载 的机制}
Transformer 训练由大 GEMM、attention、normalization、softmax、optimizer step 和通信构成。 这句话把讨论落在张量、token 和计算图状态上。读者需要知道哪些量可直接观察，哪些量只能通过日志、profile 或评估协议间接推断。
大 GEMM 更容易吃满 tensor cores；softmax、LayerNorm、dropout、optimizer 等更容易 memory-bound。 因而同一个方法在不同规模下可能改变不同的变量，尤其是sequence length、hidden size、head 数、activation size 和训练稳定性。比较方法时，必须同时给出问题规模和实验条件。
系统优化的目标是把 high-level PyTorch graph 映射成尽量少、尽量高效、尽量少搬数据的硬件操作。 这使本节具有可检验性：一个结论若不能在公式、伪代码、profile trace、benchmark 或 ablation 中表现出来，就仍然只是直觉。
因此，用硬件模型解释 LM 工作负载 应当被理解为可以实现和测试的机制，而不是一个章节标题。CS336 反复强调 from scratch，正是因为只有能被复现的机制，才可能被稳定改进。
\subsection{用硬件模型解释 LM 工作负载 的数学表达}
\[\text{tokens/s}=\frac{\text{batch}\times\text{sequence length}}{\text{step time}}\]
符号说明：tokens/s 是训练吞吐的直接指标；step time 又由 compute、memory、communication 和 synchronization 构成。
这个关系式把 用硬件模型解释 LM 工作负载 放回模型计算图。它帮助读者追踪 token、张量形状和参数规模如何进入训练或推理成本。
例如，validation loss 常被用作模型质量的 proxy，tokens/s 用作吞吐 proxy，Jaccard 或 MinHash 用作近重复 proxy，reward model score 用作偏好 proxy。proxy 不是目标本身；它只在评估协议清楚时才可引用。
\subsection{从 用硬件模型解释 LM 工作负载 到程序}
把 用硬件模型解释 LM 工作负载 写成程序后，抽象机制会变成输入、状态和输出之间的变换。下面的伪代码保留核心计算顺序。
\begin{lstlisting}[language=Python]
profile_step()
for kernel in hot_kernels:
    classify_compute_or_memory_bound(kernel)
    optimize_layout_or_fusion(kernel)
\end{lstlisting}
这段伪代码的读法，是找出 用硬件模型解释 LM 工作负载 改变了哪些状态，以及这些状态如何对应到 sequence length、hidden size、head 数、activation size 和训练稳定性。如果某个状态没有被记录，后续实验就很难解释异常结果。
\subsection{边界条件与常见误解}
\begin{itemize}
\item tokens/s 不能单独比较模型质量；更合理的是把 loss、tokens、FLOPs 和成本一起报告。
\item benchmark 要包含 warmup 和同步，否则 GPU 异步执行会误导计时。
\item 不要把小实验中的局部结论自动外推到 frontier-scale；scale、数据分布和硬件拓扑都会改变结论。
\item 不要把 benchmark 或 profile 的单个数字当作机制解释；必须同时记录实验设置、输入形状、版本和评估口径。
\end{itemize}
\subsection{与本章其他概念的关系}
用硬件模型解释 LM 工作负载 与本章其他部分的关系是互补的：术语表给出对象名称，公式给出最小数学结构，伪代码给出执行顺序，例题则把这些内容压缩成一个可检查的计算。
\subsection{怎样检验 用硬件模型解释 LM 工作负载}
检验 用硬件模型解释 LM 工作负载 时，先把抽象说法落到张量、token 和计算图状态上。与 GEMM、LayerNorm、softmax、optimizer、tokens/s 相关的对象必须能被构造、打印、计数或评分；否则实验只能得到描述性观察，不能得到可复现结论。
一个合适的小实验应当保留本节机制，同时让 sequence length、hidden size、head 数、activation size 和训练稳定性 中至少一个量发生可解释变化。输入规模不必逼真；关键是读者能手算或打印关键中间量，并说明哪个变量触发了结果变化。
随后把公式说明转成可记录的数量关系。符号说明：tokens/s 是训练吞吐的直接指标；step time 又由 compute、memory、communication 和 synchronization 构成。 如果数量只能间接观测，就必须说明使用了什么 proxy；如果使用 benchmark 或 reward，也要说明协议和随机性。
对照实验只应改变一个主要因素，例如 batch size、sequence length、vocabulary size、attention heads、filter threshold、learning rate、decoding 参数或 verifier。多个因素同时变化时，实验最多说明系统表现不同，不能说明是哪一个机制在起作用。
最后要主动寻找失败条件。教材中的成功路径容易记住，但工程中真正决定方法边界的是失败案例。本节至少要记住这些限制：tokens/s 不能单独比较模型质量；更合理的是把 loss、tokens、FLOPs 和成本一起报告；benchmark 要包含 warmup 和同步，否则 GPU 异步执行会误导计时。复习时可以把这些限制改写成反例测试。
实验结论还要放回 GPU/TPU 硬件模型 的整体结构中。用硬件模型解释 LM 工作负载 会影响训练目标、数据选择、系统成本或评估解释中的至少一项；能说清楚这种影响，才说明读者已经从名词记忆进入机制理解。
\subsection{常见混淆：用硬件模型解释 LM 工作负载}
本节容易混淆的关键词包括 GEMM、LayerNorm、softmax、optimizer、tokens/s。它们在课程材料中可能连续出现，但层级并不相同：有的命名对象，有的描述操作，有的只是指标或约束。
其中，normalization（归一化）影响优化稳定性和 residual path。
区分这些词时，可以先问它们是否改变 sequence length、hidden size、head 数、activation size 和训练稳定性。如果一个词不能对应到变量、状态、代码路径或评估协议，它在本节中就还没有形成可操作含义。
因此，复习 用硬件模型解释 LM 工作负载 时不应只抄关键词，而应写出至少一个最小例子：输入是什么，哪一步使用这个词，输出或指标怎样变化，失败条件在哪里出现。
\subsection{用硬件模型解释 LM 工作负载 的小结}
\begin{itemize}
\item 讨论 用硬件模型解释 LM 工作负载 时，应同时报告问题规模、实验设置和主要观测变量，尤其是 sequence length、hidden size、head 数、activation size 和训练稳定性。
\item 如果一个方法声称降低主要成本，要追问成本是否转移到了显存、通信、数据质量、训练稳定性或评估复杂度。
\item 公式和伪代码提供推理骨架；结论仍需要 profiler、ablation 或 held-out evaluation 支持。
\end{itemize}
\section{综合讨论：从局部机制到完整系统}
本章的主题是 GPU/TPU 硬件模型，但它不应被读成几个相邻知识点的拼接。硬件层级：从 HBM 到 registers、Warp、thread block 与 occupancy、用硬件模型解释 LM 工作负载 共同回答同一个问题：语言模型系统中哪些对象可以被定义，哪些成本可以被估算，哪些结论必须通过实验或评估来确认。
这些对象包括 GPU（图形处理器）、TPU（张量处理器）、SM（streaming multiprocessor）、HBM（high-bandwidth memory）、tensor core（张量核心）、occupancy（占用率）。读者可以把它们看成一组接口：有的接口连接文本和 token，有的连接模型结构和张量计算，有的连接数据样本和训练目标，有的连接离线评估和在线服务。接口一旦改变，后面的成本、错误类型和优化空间也会随之改变。
\subsection{硬件层级：从 HBM 到 registers 在系统中的位置}
硬件层级：从 HBM 到 registers 首先提供一个局部解释：视频宣布进入 systems portion，目标是理解 GPU、parallelization 和 inference，而不是把硬件当黑箱。 这类解释的价值在于把直觉压缩成可操作对象。一旦对象明确，读者就可以追问它的输入、输出、中间状态和失败条件，而不是只记住方法名。
其次，硬件层级：从 HBM 到 registers 会改变至少一个资源或评估变量。它可能改变 sequence length、activation size、memory bandwidth、communication volume、data mixture、reward distribution 或 benchmark protocol。这些变量在小例子里可以手算，在真实训练和服务系统里则需要 profiling、logging 和 held-out evaluation。
最后，硬件层级：从 HBM 到 registers 不能脱离边界条件理解。一个关键 caveat 是：硬件规格表给的是上限，不是可直接达到的性能。。把 caveat 写出来不是为了削弱结论，而是为了说明结论在哪些输入、硬件、数据和评估条件下才成立。
\subsection{Warp、thread block 与 occupancy 在系统中的位置}
Warp、thread block 与 occupancy 首先提供一个局部解释：CUDA 编程模型通过 thread、warp、thread block、grid 抽象硬件并行性。 这类解释的价值在于把直觉压缩成可操作对象。一旦对象明确，读者就可以追问它的输入、输出、中间状态和失败条件，而不是只记住方法名。
其次，Warp、thread block 与 occupancy 会改变至少一个资源或评估变量。它可能改变 sequence length、activation size、memory bandwidth、communication volume、data mixture、reward distribution 或 benchmark protocol。这些变量在小例子里可以手算，在真实训练和服务系统里则需要 profiling、logging 和 held-out evaluation。
最后，Warp、thread block 与 occupancy 不能脱离边界条件理解。一个关键 caveat 是：高 occupancy 无法弥补非合并访存或严重 bank conflicts。。把 caveat 写出来不是为了削弱结论，而是为了说明结论在哪些输入、硬件、数据和评估条件下才成立。
\subsection{用硬件模型解释 LM 工作负载 在系统中的位置}
用硬件模型解释 LM 工作负载 首先提供一个局部解释：Transformer 训练由大 GEMM、attention、normalization、softmax、optimizer step 和通信构成。 这类解释的价值在于把直觉压缩成可操作对象。一旦对象明确，读者就可以追问它的输入、输出、中间状态和失败条件，而不是只记住方法名。
其次，用硬件模型解释 LM 工作负载 会改变至少一个资源或评估变量。它可能改变 sequence length、activation size、memory bandwidth、communication volume、data mixture、reward distribution 或 benchmark protocol。这些变量在小例子里可以手算，在真实训练和服务系统里则需要 profiling、logging 和 held-out evaluation。
最后，用硬件模型解释 LM 工作负载 不能脱离边界条件理解。一个关键 caveat 是：tokens/s 不能单独比较模型质量；更合理的是把 loss、tokens、FLOPs 和成本一起报告。。把 caveat 写出来不是为了削弱结论，而是为了说明结论在哪些输入、硬件、数据和评估条件下才成立。
\subsection{把本章知识用于读论文和复现实验}
读相关论文时，可以把本章变成一个检查表。第一，论文是否清楚定义了输入规模、模型规模、数据来源和评估协议；第二，论文声称的改进落在哪个变量上；第三，作者是否报告了会让结论失效的设置；第四，实验是否能分离模型结构、数据质量、优化设置和系统实现的影响。
复现实验时，也应避免直接追求大规模。先用最小程序复现一个公式、一个伪代码分支或一个 profile 现象，再扩大输入规模。这样做的原因很简单：如果小规模程序无法解释中间量，大规模训练只会把错误隐藏在日志和随机波动里。
把这些检查合在一起，本章给出的不是一个固定答案，而是一种阅读 CS336 的方法：每个模型机制都要落到变量，每个变量都要落到测量，每个测量都要说明协议，每个协议都要承认边界。
因此，本章中的任何一句结论都可以被改写成一个可引用的完整句式：在给定数据、模型、硬件和评估协议下，某个机制改变了某个变量，并通过某个观测指标表现出来。这样的句子比单独背诵术语更长，但它包含了自学、复现和引用时真正需要的信息。
把本章用于自学时，可以按三遍阅读。第一遍只追踪对象和定义，确认每个术语到底指向文本、token、张量、样本、请求还是评估记录。第二遍追踪公式和伪代码，确认每一步改变了哪些状态。第三遍追踪实验和 caveat，确认哪些结论只在特定规模、数据或硬件条件下成立。
把本章用于复习时，则应反过来操作：先从一个失败案例或异常指标出发，再回到相应机制。loss 曲线异常可能指向数据、优化或模型容量；吞吐异常可能指向 kernel、通信或调度；benchmark 异常可能指向 contamination、prompt protocol 或采样设置。能够从异常反推机制，是掌握本章的实用标准。
\section{例题：把概念变成可计算检查}
\subsection{判断一个算子是 compute-bound 还是 memory-bound}
对一个 elementwise activation 和一个大矩阵乘分别估算 arithmetic intensity。
\begin{itemize}
\item elementwise activation 每个元素只做少量 FLOPs，却要读写至少一次 HBM，通常 memory-bound。
\item 大 GEMM 会重复使用 tile 中的数据，FLOPs/byte 高，通常更接近 compute-bound。
\item 优化前先放到 roofline 模型里，避免把时间浪费在错误瓶颈上。
\end{itemize}
\textbf{结论。} GPU 优化不是“多用并行”四个字，而是让数据移动和计算形状匹配硬件。
这个例题的目的不是替代完整工程实现，而是给出一种最小推理形式：把问题转成变量，写出近似公式或伪代码，再说明近似何时失效。
\section{实验设计：从小例子到可复现结论}
学习 GPU/TPU 硬件模型 不能停在概念层。一个教材化结论至少需要四件事：可构造的输入、可观测的中间量、可比较的指标，以及清楚的失败条件。下面的实验设计不是要求训练大模型，而是要求读者用小规模程序复现本章机制。
\subsection{最小输入}
先构造只触发本章核心机制的小输入。例如 tokenizer 章节可以用几个包含 rare words 和 Unicode 字符的字符串；GPU 章节可以用一个 elementwise kernel 和一个 GEMM；data 章节可以用十几篇近重复文档。小输入的价值在于它让错误可见。
\subsection{资源账本}
对同一输入写下 FLOPs、bytes moved、显存峰值、通信量或样本数量的估算。估算不需要一开始就精确，但必须说明单位和近似假设。随后用 profiler、benchmark、validation loss 或人工检查去验证估算是否同量级。
\subsection{单变量消融}
一次只改变一个因素，例如 vocabulary size、head 数、batch size、filter threshold、reward scale、decoding temperature 或 GPU 数量。若多个变量同时变化，实验只能说明系统变了，不能说明机制是什么。
\subsection{失败条件}
最后要刻意触发失败：极端 context length、错误 dtype、过小 batch、污染 benchmark、低质量数据或不可靠 verifier。失败条件常常比成功曲线更能说明一个方法的边界。
\subsection{记录与报告}
实验记录应能让另一个读者复现同一结论。最少要写清楚输入构造、代码版本、随机种子、硬件、batch shape、数据版本、评价指标和异常处理方式。报告结论时不要只写“更快”或“更好”，而要说明快在哪里、好在哪个指标上，以及是否牺牲了内存、稳定性、数据质量或可解释性。
常见报告错误是把局部现象写成普遍规律。例如一次小模型实验里的 loss 改善，不等于更大模型或不同数据分布上也会改善；一次 kernel benchmark 的吞吐提升，不等于端到端 serving latency 一定降低。教材中的实验报告应始终把结论限定在可复现条件内。
\subsection{从小实验到大结论}
小实验的作用不是替代课程中的完整训练或系统实现，而是建立一条可检查的推理链。先在小输入上确认变量方向，再在中等规模上确认数量级，最后才讨论是否能外推到更大模型、更多 token 或更复杂 workload。缺少中间层级时，大结论通常只是在复述直觉。
因此，实验报告最好同时包含正例和反例：正例说明机制在受控条件下怎样工作，反例说明条件稍变时哪里失效。对 CS336 这样的课程，反例尤其重要，因为 language model 的错误常来自多个层面叠加：数据分布、优化稳定性、硬件瓶颈、评估协议和服务负载都可能同时改变观察结果。
还要记录资料版本与复现边界。公开视频、课程页、配套代码和 PDF 可能在不同时间更新；复现实验应写明使用的材料版本，并说明哪些结论来自课程事实，哪些是为了帮助理解而加入的解释性推导。这样才能把学习笔记变成可检查的技术记录。
若复现实验与课程结论不一致，首先检查输入、版本和评估协议，而不是立即修改模型假设。
\section{课程资料与回看路径}
教材正文已经把主要概念、公式和实现逻辑组织成连续叙述；本节只提供回到课程材料的阅读路径。读者复习时可以先完成前文练习，再按下面的时间段或资料组核对细节。
\subsection{视频回看路径}
\begin{longtable}{p{0.20\linewidth}p{0.46\linewidth}p{0.26\linewidth}}
\toprule
\textbf{时间} & \textbf{关键词} & \textbf{复习任务}\\
\midrule
00:00:05-00:06:39 & going, think, your, faster, kind, systems, understand & 回看定义、例子、推导或 caveat 在该段中的位置。\\
00:20:19-00:26:48 & matrix, that's, multiply, it's, going, very, memory & 回看定义、例子、推导或 caveat 在该段中的位置。\\
00:39:39-00:46:02 & going, scaling, think, quantize, matrix, it's, factor & 回看定义、例子、推导或 caveat 在该段中的位置。\\
00:52:04-00:58:54 & memory, read, going, threads, matrix, think, it's & 回看定义、例子、推导或 caveat 在该段中的位置。\\
01:11:58-01:18:25 & going, then, think, attention, compute, very, softmax & 回看定义、例子、推导或 caveat 在该段中的位置。\\
\bottomrule
\end{longtable}
这些时间段用于定位视频中的讲解顺序。严肃引用时，应回到视频片段确认讲者表述、板书或幻灯片内容。
回看视频时，不必逐字抄录字幕。更有效的方法是把每个时间段判定为定义、例子、推导、代码解释或 caveat，再把它放回本章对应小节。
\subsection{课程资料阅读路径}
\begin{longtable}{p{0.14\linewidth}p{0.48\linewidth}p{0.30\linewidth}}
\toprule
\textbf{资料组} & \textbf{标题/页块} & \textbf{关键词}\\
\midrule
official\_01 & Lecture 5 GPUS CS336 Tatsu H & memory, gpus, scaling, threads, fast, compute, https\\
official\_02 & Recap: GPUs – what are they and how do they work ❖ GPUs are massively parallel – same ... & memory, gpus, fast, same, make, workloads, control\\
official\_03 & Trick 1: Low precision computation If you have fewer bits, you have fewer bits to move & memory, https, precision, access, operations, trick, read\\
official\_04 & Tiling math Non-tiled matrix multiply: each input is read 𝑁 times from global memory ... & tiling, matrix, memory, softmax, part, attention, computation\\
\bottomrule
\end{longtable}
课程资料通常给出更稳定的标题、公式、图或代码结构；视频则补充动机和口头解释。两者合起来构成本章的事实边界。
阅读资料时，可以把每个资料组拆成三个对象：一个定义，一个公式或伪代码动作，一个边界条件。这样比逐字翻译页面或脚本更容易形成可用知识。
\section{实践说明、历史位置与风险}
GPU/TPU 硬件模型 在 CS336 的课程结构中连接基础机制和规模化问题。基础机制解释方法为什么成立；规模化问题检验它在更大模型、更长上下文、更复杂数据或真实 serving workload 下是否仍成立。
历史位置不等于模型年表。更重要的是识别问题如何迁移：早期方法解决小规模建模问题，现代方法常常解决同一问题在大规模数据、GPU 集群、长上下文或人类偏好约束下的新形态。
\begin{itemize}
\item 资料版本限制（version risk）：课程页、公开视频和配套资料可能在学期中更新；引用时应注明访问日期或资料版本。
\item 测量口径限制（measurement risk）：FLOPs、tokens/s、benchmark score、reward 等数字只有在协议完整时可比较。
\item 规模外推限制（scaling risk）：小模型或小 batch 的机制不必然外推到 frontier-scale。
\item 数据分布限制（data risk）：数据来源、过滤和去重策略会改变模型行为，也会改变 evaluation contamination 风险。
\item 系统风险（systems risk）：性能剖析结果依赖硬件、driver、kernel、shape 和 batch distribution；脱离 workload 的性能数字不可直接复用。
\end{itemize}
\section{复习要点}
下面的问题把本章内容压缩为可反复检查的条目。每个条目都应能回到前文的解释、公式、伪代码或例题。
\begin{longtable}{p{0.16\linewidth}p{0.38\linewidth}p{0.38\linewidth}}
\toprule
\textbf{条目} & \textbf{对象} & \textbf{检查动作}\\
\midrule
条目 1 & GPU（图形处理器） & 定义它；写出一个最小例子；说明一个 failure mode；指出它影响哪个资源或评估账本。\\
条目 2 & TPU（张量处理器） & 定义它；写出一个最小例子；说明一个 failure mode；指出它影响哪个资源或评估账本。\\
条目 3 & SM（streaming multiprocessor） & 定义它；写出一个最小例子；说明一个 failure mode；指出它影响哪个资源或评估账本。\\
条目 4 & HBM（high-bandwidth memory） & 定义它；写出一个最小例子；说明一个 failure mode；指出它影响哪个资源或评估账本。\\
条目 5 & tensor core（张量核心） & 定义它；写出一个最小例子；说明一个 failure mode；指出它影响哪个资源或评估账本。\\
条目 6 & occupancy（占用率） & 定义它；写出一个最小例子；说明一个 failure mode；指出它影响哪个资源或评估账本。\\
\bottomrule
\end{longtable}
\subsection{综合自测}
\begin{itemize}
\item 我能否不用课程原句，独立解释本章中心问题？
\item 我能否把本章至少一个公式改写成可运行的估算函数？
\item 我能否指出一个看似改进但实际转移成本的方法？
\item 我能否回到视频和课程资料，找到支持本章关键论断的位置？
\item 我能否设计一个最小 ablation，验证本章一个 caveat？
\end{itemize}
\section{总结与延伸}
\subsection{本章小结}
\begin{itemize}
\item GPU/TPU 硬件模型 的核心不是记忆术语，而是理解对象、变量和机制之间的关系。
\item 视频给出讲解顺序，课程资料提供可核对的公式、图、代码或页面结构。
\item 复习时应把每个术语连接到至少一个变量、一个实现动作和一个失败模式。
\end{itemize}
\subsection{延伸解释}
本章中的延伸解释用于连接课程材料中的概念，例如说明公式里的 proxy、解释系统 tradeoff，或把多个材料片段串成一个完整推理链。若需要严肃引用，应回到课程视频和配套资料核对原始表述。
\subsection{练习}
\begin{itemize}
\item 定义题：用中英双语解释本章任意五个重要术语，并说明它们属于 compute、memory、data、optimization、evaluation 还是 deployment 账本。
\item 推导题：选择本章一个公式，逐个解释符号、单位和近似假设，并说明哪个变量最难在真实系统中测量。
\item 实现题：把本章伪代码改写成一个最小 Python sanity check，不要求训练大模型，但要输出可检查的中间量。
\item 资料题：从“课程资料与回看路径”中选择一个视频时间段，回到原始视频核对讲者例子，并记录是否存在字幕误差。
\item 评估题：为本章方法设计一个 ablation 或 benchmark，说明控制变量、数据版本、指标和 failure criteria。
\item 边界题：说明本章一个结论在更长 context、更大 batch、更多 GPU 或更复杂数据分布下可能如何失效。
\item 综合题：把本章内容和前一章或后一章连接起来，写出一个端到端训练/推理 pipeline 中的依赖关系。
\item 批判题：指出一个容易被营销材料夸大的说法，并用本章的资源账本或评估协议反驳它。
\end{itemize}
\end{document}
